\chapter{Computational philosophy}

We, now at least have a rough idea on where and what to look at, at least in terms of defining intelligence, what is the general beholden conception of what is \underline{considered}, in antique terms, \textit{artificial intelligence}. Now we need to put that constructions in good use. It will also reinforce our understanding, and our formulation of the \textbf{actual} implementation for the actual prototype and framework of which constitute artificial intelligence. Remember, for myself and ones reading this - our goal is to put this as a documentation toward creation, and not crackpot theory and arguments. This is imperative. 

The first `good use' that will be touched upon in this section, is the method, notion and usage of historically a very familiar concept as of date, in which many inventions, many works and rather performance scaling works - the computational process, the \textit{computation} concept itself, the concept of a \textit{simulation}, and the surrounding \textit{computer science} of things. Those, in our model, for now, comprise sufficiently in the \textbf{existential facilities} mentioned, at least of the current theory, since it fuels the creation and adaptation of the learning machine, and the adaptive machine we theorized somewhat from what was available. However, the philosophy, fundamental arguments, foundational principles, and overall guarantees, as well as the epistemology of such is not digged upon in much of its existence. There are too many dark clouds need to be settled, while less than that is the clear vision of what can be deemed important to consider if computation is the universal imitation or not. Note that, however, the topic on which we called \textbf{simulation} is much more difficult to handle, as it is larger than just computer simulation. You do not need a specifically advanced computer to somewhat put upfront a general interpretation of the relativity of the sun and planets - you need a blanket and two balls. Hence, in the context we deliver it, the term would be dubbed as \textbf{computer simulation}. 

Overall, what is \textit{computation}? What is meant by being computed? Why we need it, why we want it (for our constructs), what gives, what takes, and what is the essence of it? More importantly, can we \textbf{justify} its capability and range of expression - as in, how much complex and details it can get? And most of such, does the theory of neural computation, or \textit{intelligence via computable machines} theory makes sense, if there exists a justification that the eventual outcome is indeed computational? We will attempt to raise those questions and more, and try to put forth some arguments and answers to it. As of the old quote of myself when encountering this topic. 
\begin{quote}
    It is customary to think of the machine, of the current time as \textit{computers}, specifically, and most of the time directed to, the \textit{binary-based computer}. However, this is a loose generalization, accompanied by many, if not a lot of assumptions, theories and philosophical presumptions that effectively guarantee specific effectiveness at hand for such computational system. It is thereby imperative to contain an expedition toward computation as a whole, examining the philosophical theory of computation, and its consequently computational ideas. We would also want to adequately familiar ourselves with particular philosophies where mathematics step in, and try to ground the thoughts. 
\end{quote}
\section{The question of computation}
Computation in history is used almost synonymous to calculation. Now it is the pillar on which computers are defined upon. In fact, in [Jan Van Leeuwen, 2015], it is stated that: 
\begin{displayquote}
    Computation is traditionally about calculating functions and solutions to equations, with or without the aid of a calculating machine. The mathematical interest for computation changed character as the result of Hilbert's famous question [13] whether there exists a finite decision procedure for first-order logic. While solving this problem, Turing [22] developed his groundbreaking model of computation that has shaped our thinking of computation since.
\end{displayquote}

It is also acknowledged that we in fact, do not have a general idea on what can be considered computation by definition, by functionality, or else, simply because the effort, or at least the insight to such problem, is relatively scarce. This issue is problematic, but pretty much interesting to look at - because at the end, our goal is artificial construct on machine, and the most powerful machine of date is the, well, \textit{computer}. If we are not even sure what is computation, can we even be sure that the current computational machine can sustain our construct? 

\begin{displayquote}[Ian Horswill, 2008]
    Computation isn't tied to numbers, acronyms, punctuation, or syntax.  But one of the things that makes it so interesting is that, in all honesty, it's not entirely clear what computation really is.   We all generally agree that when someone balances their checkbook, they're doing computation.  But many people argue that the brain is fundamentally a computer.  If that's true, does it mean that (all) thought is computation?  Or that all computation is thought?  We know that the “virtual” environments in computer games are really just computational simulations.   But it's also been argued that the real universe is effectively a computer, that quantum physics is really about information, and that matter and energy are really just abstractions built from information.  But if the universe is “just” computation, then what isn't computation?  And what does it even mean to say that something is or isn't computation?
\end{displayquote}

So, just what is computation? And can we approach from the top-bottom to what is considered by the time Turing and Church established their computational theory?

\section{Computation at the first level}
First and foremost - the definition of \textbf{calculation} is limited, at least in our terms. It strictly means the operation of elementary computation - the treatment of evaluating the language of descriptions, mathematics, specifically, with quantified notion of \textbf{numbers}. Essentially, if we think of the hierarchical structure of the abstract quantification language, we can say calculation is the barebone of operation conductible on such abstract language system. 
\subsection{Calculation}
The concept of \textbf{calculation} is ancient. It branches off, at least for what is considered the historical record of ancient civilization utilizing the tools of mathematics, as a notion manipulating the quantified - with the intended result to bring forth the presentable, \textit{existence} quantification of such. In other way of thoughts, calculation is the \textit{first} ever abstraction options available for our development of mathematical treatments in the early age, with the ability to manage quantifications, lies already in the first ever abstraction itself. For mathematics as the science of patterns, abstractions, and generally in different perspective, a reverse engineering of reality to be interpreted by human, calculation is the companion of the first layer of abstraction, for one is infinitely sufficient for others layer of the Tower of Babel \footnote{The analogy here is kind of funny, but it makes quite the sense - even if you build everything on top of those layers, the foremost primitive operation is still applicable and is of importance to the operation of the above layer. In fact, there are some core conceptual operations sets that is prominent in almost all layer perceivable - similar to how arithmetic has addition, and vector space still have addition.} on top of it. Sometimes we can even interpret calculation as a primitive type of \textit{simulation}, per se, but that is not in focus of this section (we still, will discuss it though). 

We can construct its foundational setting of calculation, and it is called \textbf{arithmetics}, as per recognition of the historical development of mathematics. By the two-fold specification, we define such as followed:

\specification{Arithmetics}{Representational alphabet $A[n]$ of size $n$, and objects $O_{1},\dots$ of quantification on such alphabet. This can be thought as the base symbol, for base-10 (modern), for example, the quantification on $\mathbb{N}$ contains 10 symbols.}{Defines operations on object of the alphabet such that it is well-defined (operations conducted on the alphabet stays in the alphabet), constructive (objects using the alphabet can be created using those operations - this is reflected, for example, in the choice of either base-2, or base-10, or base-16, and for certain construction as Peano's axioms), internal interaction rules, and well-behaved.}

From the definition of arithmetic, we construct the notion of calculation:

\begin{definition}[Calculation]
    We refer to calculation as the \textit{exact operations} of arithmetics. For exact alphabet, exact objects, and exact rules, deduce and apply operations by specific configuration and means to produce a new object lies in the same alphabet, and with well-defined behaviours, including \textbf{undetermined cases}. 
\end{definition}

From calculation, you can then define \textbf{relations} on them, or at least properties operators like $<,>,\geq,\dots$ and else, but the current layer does not have that capability yet. We need something \textit{stronger}. However, we might want to clarify that we use, at least in this level, certain form of correctness of logic circuit - the realization of logic in operational form. Between the logic and their phrases of truth value, there exists the operator $\mathrm{Eval}(P)$ in which evaluates the expression $P$, for it to be constructed from those property operators. However, since our machine is much weaker for those operations, the singular operator that is applicable for evaluation, is the equivalent relation $=$. Then $\mathrm{Eval}_{=}(P)$ returns if the relation is true for the expression $P$ given by the alphabet with the value $T$ or $F$. Because we operate on a machine model which would be assumed to be error-free, the operations work spectacularly, $\mathrm{Eval}(P)=T$ for all cases. This logic aspect is outside the layer itself, for it to evaluate based on an entirely different alphabet, rules, and evaluations, and is much more complex than the current machine can permit. Hence, this addition only works in favour of completeness. We will show you how the more `advanced' layer can create its own logic module. Then the first strong assumption here, would be that there are no expression  

Why we even bother with this kind of `logical module correctness' anyway? Well, it is because we are actively trying to make a model much more limited than the theoretical aspect there are. If you work with the language of arithmetic, you can see that our construction do allow for certain statements like $1+6=59$, for example, simply because there exists no guarantee about the correctness of what is manipulated by the layer. The abstract mathematic resolves this issue, however, by using logic, which brands $T,F$ to the statement. Because we want to construct theoretical layers that would then actively be used to theorize on the actual \textit{operational machine}, we would not want to deal with, right at the very bottom of the operation, the fact that somehow we cannot rule out the existence of false expression - and for operating machine, the `logic' module is not expressible at this level, which is troublesome. Hence, from a perspective, this is needed for justification of trueness. 
\section{Inspecting computation}

\subsection{Quick history}

Per review of the concept of computation, the specific definition that we aim for is utimately decided and define by either Godel (1934), Church (1936), Post (1936) and A. Turing (1937) by defining the conceptual machine that do computation, rather than the concept of computation itself in some case, and the evaluation required to describe the concept in others. Overall: 
\begin{itemize}
    \item Godel defined it in terms of the evaluation of recursive functions. 
    \item Church defined it in terms of the evaluations of lambda expressions. 
    \item Post defined it as series of strings sucessively rewritten according to a given rule set. 
    \item Turing defined it the sequence of states of an abstract machine with a control unit and a tape (the Turing machine).  
\end{itemize}
In the time that these men wrote, the term meant the mechanical steps followed to evaluate mathematical functions. Removing the mechanical constraint to the abstract level, it means the step of a machine to evaluate mathematical functions by certain rules. In such sence, computation hasn't changed since then. 

Normally, the scheme of searching for computation is not to define, or to discover the essence, at least figuratively, about computation. Rather, it is to create an actor that we are able to define its operations as computation, and itself a \textit{computing machine}. This is most powerfully expressed using a Turing machine construct. Formally, the original construction of a Turing machine (automata) is as followed: 
\begin{definition}[TM]
    A \textbf{Turing machine} $M$ is a 6-tuple $\{Q,\Sigma,\Gamma, \delta, q_0, H\}$ consists of:
    \vspace{2mm}
    \setlist{nolistsep}
    \begin{enumerate}[noitemsep]
        \item $Q$: Finite set of states. 
        \item $\Sigma$: Finite set of input alphabet, excluding $\vartriangleright, \square, \rightarrow, \leftarrow$. 
        \item $\Gamma$: Finite set of tape alphabet. $\Sigma \cup \{\vartriangleright, \square\}\subseteq \Gamma$. $\Gamma$ still excludes $\rightarrow, \leftarrow$. 
        \item A transition function $\delta: (Q-H)\times \Gamma \to Q\times (\Gamma \cup \{\rightarrow, \leftarrow\})$ such that the tape head never falls off or erases the $\vartriangleright$, where $\vartriangleright$ is the computational time.
        \item $q_0$: The \textbf{start state} (belong to $Q$). 
        \item $H=\{q_{acc}. q_{rej}\}$: The set of \textbf{halting states}, subset of $Q$. 
    \end{enumerate}
\end{definition}

This system, at least for Turing as above, operates on the principle of first action. A \textbf{configuration} or \textbf{global state} of a Turing machine is represented by a triple $(q,x,i)\in Q\times \Gamma^{*}\times \mathbb{N}$, indicating that the machine's internal state is $q$, the tape contains the string $x$ followed by an infinite sequence of blanks, and the head is located at position $i$. 

%\begin{interject}[TM, II]
    %The Turing machine (TM) perhaps can be referred to as a \textbf{linear machine}. To be fair, a linear machine is specifically easier to resolve than a non-linear one. However, such non-linearity is quite important, certainly, if the problem scales up, or rather, when the scope of the problems, as well as the process grew. Approximation helps, certainly, but there are restrictions. What is the impulse for Turing machine to be a linear machine? Are there reasons of such, for example, the perservation of the totality of the machine itself, or that the behaviour (for example, Turing-Church hypothesis) would be violated. Rather then, what can be said of \textbf{linearity} and \textbf{nonlinearity} for such TM, and then finite state automata at large? Is this linearity is resolved only in its operating space, or even the operating principle and class of response? 
%\end{interject}

\begin{displayquote}[Peter J. Denning]
    The standard formal definition of computation, repeated in all the major textbooks, derives from these early ideas. Computation is defined as the execution sequences of halting Turing machines (or their equivalents). An execution sequence is the sequence of total configurations of the machine, including states of memory and control unit. The restriction to halting machines is there because algorithms were intended to implement functions: a nonterminating execution sequence would correspond to an algorithm trying to compute an undefined value.
\end{displayquote}

Throughout time, computation has been looked at as calculation, symbol manipulation, information processing, and then as a process, often in connection to a model of computation that is considered as being realistic now or in the future. It evolves from when it is considered to be mental works that include numbers - and as being mental operations, it can only be conducted by human, to now when the role is switched to creating computers that compute, not involving and mental task of all. 

However, for representative meaning, let's have a look at how current literature views the concept. In \textit{Introduction to Theoretical Computer Science (2023)} by Boaz Barak, computation is defined by:

\begin{displayquote}
    To a first approximation, computation is a process that maps an input to an output. 
    
    When discussing computation, it is essential to separate the question of what is the task we need to perform (i.e., the specification) from the question of how we achieve this task (i.e., the \textit{implementation}). For example, as we've seen, there is more than one way to achieve the computational task of computing the product of two integers.
\end{displayquote}

Afterward, the main section is reserved for algorithm being prominent as such. We can take this, and for Denning's mention of Djikstra's clarification between computation and algorithm that computation is the \textit{realization of algorithm} in one sense or another. So how do this work? 

The truth is, this is harder than it looks. Typically, the question of computation, algorithm and else resides by the definition within the mechanical computing machine, or the abstract computing machine, for example, Turing machine (TM) or finite automata construct. We determines four subjects of interest as we see it such. By far, algorithm, computation, models, and mathematics. Out of this, mathematic is the language of abstracted quantization. By abstract, we mean the absence of physical realization or existence. A \textbf{model} can be mathematical or not, but inherently it follows certain logic - realistically though, the definition of logic ultimately can be interpreted to the deductive reasoning of human - but we would like to use the simplification of such, then a \textit{model} is a realization or abstract inference of certain construct of interest. In such definition, mathematical model is the abstract realization for inferencing certain construct of interest, via the language of mathematical quantization. The interaction and governing realization of the model is described, using \textit{algorithms}, which is to serve specific process in the model itself. And to implement all of them, computation is the realization to quantized reality and case, even with the absence of physical realization. This will be reflected in our structure. 

\subsection{Modern history (1900s --- Current) - Computational theory of mind}

Such framework, so powerful it is made way for a new kind of hypothesis. Could a machine think? Could the mind itself be a thinking machine? The computer revolution transformed discussion of these questions, offering our best prospects yet for machines that emulate reasoning, decision-making, problems solving, perception, linguistic comprehension, and other mental processes. Advances in computing raise the prospect that the mind itself is a computational system—a position known as the computational theory of mind (CTM).

The classical computational theory of mind (CCTM) \cite{compmodel2016} relies on the foundational concept of Turing machine. In this sense, it is noticed that a Turing machine contains a \textit{central processor} and a \textit{memory store}. A Turing machine manipulates primitive symbols drawn from a finite alphabet. The central processor inscribes or erases symbols in memory locations, depending on the current contents of a memory location and the processor's own state. A \textit{machine table} encodes precise rules dictating how the central processor should proceed. 

According to CCTM, the mind is a \textit{computational system} that resembles the Turing machine in important aspects, which means the formalism is not followed, yet the central theme and ideas are implemented for the prototypical structure of the proponents' model. By CTM, additionally, it is also accepted that the mind computes over symbols, or relatively physic'less' state. For such, it is then considered that primitive symbols can be compounded into complex symbols, in ways or procedures that resemble either complex construction in predicate logic or first-order formalism, or any system with propositional structures embedded. This idea, famous by Jerry Fodor, is called \textit{Mentalese}, though it is associated more with the representational theory of mind. Thereby, it is reasonable that the mind can be described by a computing system, and not a \textit{computer} - the latter follows that subject is programmable, but philosophically it is not required to be programmable in explaining the mind. 

This way of view, however, seems to be too restrictive. This is, by certain counts [Stanford, M Rescoria, 2024]: 
\begin{itemize}
    \item Turing machines execute pure symbolic computation. The inputs and outputs are symbols inscribed in memory locations. In contrast, the mind receives sensory input (e.g., retinal stimulations) and produces motor output (e.g., muscle activations). A complete theory must describe how mental computation interfaces with sensory inputs and motor outputs.
    \item A Turing machine has infinite discrete memory capacity. Ordinary biological systems have finite memory capacity. A plausible psychological model must replace the infinite memory store with a large but finite memory store.
    \item Modern computers have random access memory: addressable memory locations that the central processor can directly access. Turing machine memory is not addressable. The central processor can access a location only by sequentially accessing intermediate locations. Computation without addressable memory is hopelessly inefficient. For that reason, C.R. Gallistel and Adam King (2009) argue that addressable memory gives a better model of the mind than non-addressable memory.
    \item A Turing machine has a central processor that operates serially, executing one instruction at a time. Other computational formalisms relax this assumption, allowing multiple processing units that operate in parallel. Classical computationalists can allow parallel computations (Fodor and Pylyshyn 1988; Gallistel and King 2009: 174). See Gandy (1980) and Sieg (2009) for general mathematical treatments that encompass both serial and parallel computation.
    \item Turing computation is deterministic: total computational state determines subsequent computational state. One might instead allow stochastic computations. In a stochastic model, current state does not dictate a unique next state. Rather, there is a certain probability that the machine will transition from one state to another.
\end{itemize}

The view of computationalism by CTA, coupled with functionalism, in which describes the state of mind or logical construct relevant as its process than what it is endowed upon by what it is, forms \textit{computational functionalism}. In this view, lies the analogy between some features of minds and computers. [Putnam, 1960, 1967a,1967b] puts this into some details. He noticed that the state of Turing machines are individuated in terms of the way they affect and are effected by other Turing machine states, inputs and outputs. By the same token, mental states are individuated by the way they effect and affected by others. It is then reckoned that mental states can be characterized functionally, like those of Turing machine, though he added that the mind might be something "quite different and more complicated" than a Turing machine. The final idea resides in describing mental states being the \textit{probabilistic Turing machine states}. 

Overall, however, the computational theory of mind suggests that \textit{computation} is the key, and \textit{computational model/system} is the truth model of modelling the mind, or any actors of sufficient complexity to be called exhibiting artificial intelligence. Though such discussion in philosophy does not translate well to what would then be established, it is still a worthy ground of consideration. How do we translate this into being, if we are to keep this point of view? 

\subsection{Abstraction of computational processes}

In the past, computation is quite the same: just that it is regard as about calculating functions and solutions to equations, with or without the aid of a calculating machine [J. V. Leuuwen, 2015]. This is somewhat presents in a lot of fields - from the meaning and practice of computational physics, computational chemistry, to computational biology and computational astrophysics. Most of the time, it is about solving problems and a system of functions, with actual numbers. Unlike calculation, the classical definition of computation seems to directly infer to the point of it being \textit{abstract calculation} instead - with the difference between $7+4$ and $x+y$ for the latter being computation is the absence of case-specific quantification. So $x,y$ can be interchanged at times. How we even define this? We sure have one of the very early definition we will make use of for now. 

\begin{definition}[Abstraction, 1]
    The \textbf{algebraic abstraction} is defined as the \textit{abstraction} of calculation. For any numerical representation alphabet $N$ and their operators, define an obstensive definition of an operation named \textbf{procedure} such that constructions from the alphabet can be applied upon, and the result of the procedure is well-defined within the alphabet. 
\end{definition}

If this feels familiar, you would be right. This is similar to the concept of a \textbf{function} definition. The name \textit{procedure} would be more similar case-by-case with the concept of relation, however, since there will be (as we will see) for situations and specifications that would not guarantee the analogy to work. 

What is the subject of this new `computation level 1'? For those operations are now clear, we might want to see some of it in the similar structure as the above box. A procedure can be either a function, or not, by the definition of function meaning a many-to-one operation in the same computational scheme. The subject of this abstraction, is not constructions (or, quantified, specified numbers) but the \textit{collections of constructions} from a given alphabet, in which case we have the numerical alphabet as mention in calculation. For this, we need a new alphabet. Fortunately, this is conducted formally using the usual abstraction alphabet. Here's the thing: aside from the primitive layer occupied per descriptive definition being calculation, the others layers can use abstracted, non-functional language and their alphabets but still retains functionality, because, they are abstractions, which in turns, result in conventions and a top-down solution. \footnote{Per my secluded knowledge and insufficient understanding, perhaps this is why there existed a period where the foundational structure of mathematics is so focused on the objective notion of the formalization and rigorous definition of arithmetics, considering it is \textbf{unfairly, stupendously strong} system of abstraction right at the beginning.} Up to this point, we start to think a bit\dots more abstract. Now we can think of objects in question, as \textbf{sets}. We will not only take object of specific alphabet construction, but the set of \textit{all objects} either with specific evaluation or not. We also deal with \textit{variables} as the named suggest, the abstract placeholder of unspecified exactness. 

From that, came an understanding, and we might as well formulate our notion of this layer similarly. 

\specification{Abstraction, I}{Abstracted representation alphabet $\mathcal{C}[m,1]$ of size $m$ for convention on set of objects $\{C^{(1)}_{1},\dots, C^{(1)}_{k} \}$ abiding the alphabet $A[n]$ in \texttt{Calculation}. Define \textbf{variables} as a placeholder, such that specific subject in such variable $x$ abides the alphabet it is in, and \textbf{constant} as the denoted alphabetical construction that is required to be given specific meaning, without ambiguity - that is, it is an inherent exact unit, an is related directly to the previous abstraction (in this case, \texttt{Calculation} layer.). Constant can be constructed from variable by binding a specific alphabet `naming' to singular quantification.}{Defining \textbf{procedure} and its special case, \textbf{function}\footnote{I did not specify this, but for transparency or sufficient detail, we might want to re-implement the definition and actual construct of what is considered procedure and function as it is, to elude confusions. } such that it abides well-defined rules, constructive rules, function self interaction rules, and the \textit{semi-recursive} rules (all operations from the previous layer are applicable to the function object). Additionally, by the equivalent relation $=$, procedures can be reduced by operation to the lower level during exact evaluation.}

There's the oddity of this abstraction with variable, which is fairly interesing. If we are to follow this thought, then when variable are switched to exact, it becomes constant. It now has unchanged value throughout the operation until the results come about, and new action begin with either that value, or some others. We might want to say then, variable under exactness of this level, represents \textit{temporary constant} in the operation. 

Then oddly, the definition for computation is quite the same as the definition of calculation for arithmetic. 

\begin{definition}[Computation, I]
    \textbf{Computation} is the exact `implementation' of the specified procedure, under specific representational abstracted alphabet of choice, in this case, we call as \textit{non-valued alphabet} in the sense that it is not embedded with any quantification, and satisfies operation properties and well-define property.
\end{definition}

What is the end goal of computation, at this point? Even though being very, very limited, and allow no algebra, computation is similar to calculation for arithmetic - describe, or exhibit exact operation for the corresponding abstraction level. In one way or another, if abstraction's goal is to generalize, conceptualize, complex-making structures to be presented, then computation, accompanied by the abstraction layer is the realization of such abstraction. The word `exactness' appears too much in the above passages that when you think of it, computation is the exact, specific implementation such that there exists the non-ambiguous presentation at the end. The basic example, at this stage permitted, is for the function, for example, $f(x)=3x+4$ to be computed, a specific number $x$ must be substituted for specificity. The \textit{graph} of the function is created not by abstractly stating it, but by computing it enough for it to be drawn. To \textbf{compute} means to realize from the abstract state. It gains its expressibility by abstraction of its according construction layer, but its operation is dictated by specificity - a computing machine cannot work if you input it nothing, but it remains powerful enough and expressive enough to handle what it is expected. For now, this realization is limited to the direct evaluation of specific `value', to specific construction, and this is the definition for the current layer. Can we \textit{expand} it any further, though? \footnote{We also note that, at later section, we would also like to commit a fair amount of works onto the discussion of the machine in which operate those computations. In one sense, a computational model, and computation is different, and that is to be seen.}

\subsection{Inspecting abstraction}
We used a lot of notions in this section taking turns on the abstraction of the logical unit, representational unit, and the alphabetically-based operations of computation. Seemingly create an alphabet, and then a structure of arithmetics as its baseline abstraction, then we expand it into further terms and objects encapsulated the previous abstracted structure. However, what are some thoughts on the topic of abstraction that align with ours, or at least, refute some of our approaches?

\begin{displayquote}[Hartry, Field]
    Sets as well as directions and numbers are to be introduced by abstraction. E.g. simple abstraction: it is suitable for us saying that our talk of directions refers to parallelism. - But that does not quite work accordingly for numbers as it does for non-numeric talk (and "non-set theory").
\end{displayquote}

\begin{displayquote}[Frege, Stuhlmann-Laeisz]
    Abstraction is a process to identify each other in the various elements of a region, e.g. same color, same size, same shape. New objects are gained by abstraction of the partial identity. ((s) Such objects which are only about numerical equality or color identity are individuated.) If $A$ is such an abstraction, then there is by definition a (base) object $a$, so that the following is true: $A$ is the $F$ of $a$.
\end{displayquote}

\begin{displayquote}[Thiel]
    Abstraction is a purely logical process, an operation with statements, the logical character of which is revealed by the change from the structure of the complicated initial statement to the structure of the new statement. Frege understood this first. Tere are statements not only about numbers, but also about sets, functions, concepts, situations, meaning and truth value of a statement, about structures.

    We also had no reason to conceive the digits as "name" of numbers, so that 4, IV, and $||||$ would denote the same number four, as it was assumed in the traditional philosophy of numbers. 
\end{displayquote}

Aside from some treatment of mathematical construct, with $\lambda$-calculus (Church) in proponents, the direct definitions and arguments on abstraction \footnote{See \href{https://philosophy-science-humanities-controversies.com/table-concepts.php}{Dictionary of Arguments - Philosophical and Scientific Issues in Dispute, Concepts A-C table}.} seems to align with what we reasonably want to enact. Thiel also garnered sufficient groundwork for discussion abstraction as a universal process in the abstracted notion of the theoretical linguistic world that we call within such, the symbols and manipulation that is yet to be tautologically same as mundane tasks. 